%% kapa: exam=1 round=21 year=2010 period=2 form=A subject=appraisal_law subject_ko=부동산관계법규 seq=30 number=70
%% source: Q-Net 감정평가사 기출문제 (cst00309) artlSeq=5024181 file=extracted/1008/5024181/2010년 제21회 감정평가사 1차시험 기출문제/2교시.hwp
%% conversion: kordoc 4.24.51 -> tools/gen1.py (auto); review=proofread-gpt-sol; flags=none
\begin{question}{70}
건축법령상 용도별 건축물의 종류를 연결한 것으로 옳지 않은 것은?
\par\noindent\hspace{2em}\uline{\makebox[13em][l]{건축물의 종류}}\quad\uline{\makebox[15em][l]{건축물의 용도}}\par
\begin{choices}
  \item \makebox[13em][l]{다가구주택}\quad -- 단독주택
  \item \makebox[13em][l]{여객자동차터미널}\quad -- 운수시설
  \item \makebox[13em][l]{무도학원}\quad -- 위락시설
  \item \makebox[13em][l]{식물원}\quad -- 동물 및 식물 관련 시설
  \item \makebox[13em][l]{운전학원}\quad -- 자동차 관련 시설
\end{choices}
\end{question}
\answer{4}
